% Source: https://artofproblemsolving.com/community/c3h1614612_last_two_digit_problem
% Stable identity: aops-topic-1614612-q01
First $33^2=1089=1000+89$, so $33^{288}=1089^{144}=(1000+89)^{144}$. Expanding by the binomial theorem, every term except $89^{144}$ contains a factor $1000$, so $33^{288}$ and $89^{144}$ have the same last two digits.

Now $89=90-1$, so
\[
89^{144}=(90-1)^{144}=\sum_{j=0}^{144}\dbinom{144}{j}90^j(-1)^{144-j}.
\]
For $j\ge2$ the term contains $90^2=8100$, a multiple of $100$. The terms with $j=0$ and $j=1$ are $1$ and $-144\cdot90=-12960$. So
\[
89^{144}=1-12960+100K=-12959+100K
\]
for some integer $K$. A number of the form $-12959+100K$ ends in the same two digits as $-12959+13000=41$. So the last two digits of $33^{288}$ are $\boxed{41}$.
